\section{The counterexamples}
\label{sec:counterexamples}

We collect the counterexamples to the conclusion of the no-go theorems, each with the hypothesis that
fails and the reason why it must fail.

\emph{Counterexample 1: the Lorentz super-loop.} $\mathbb L_{\rm LL}$ contains the Lorentz group
$SL(2,\mathbb C)$ and eight spinor charges that carry internal quantum numbers---isospin $\tfrac12$
and electric charge $\pm\tfrac12$---and Lorentz quantum numbers---spin $\tfrac12$, one chirality,
spinor number $-1$. The bracket of two spinor charges is a Lorentz generator, Eq.~(\ref{eq:ss}), so the
structure is not a direct product of a Lorentz part and an internal part. For a Lie algebra this is
impossible (Theorem L); for a non-associative structure it happens, and uniquely (Theorem M(b)).
Hypotheses that fail: H-grp (the loop is not associative), H-op and H-Lie (Theorem R: there is no
realization by operators, and the Jacobi identity fails, Eq.~(\ref{eq:jac6})).

\emph{Counterexample 2: inversion as an element.} The unit $\mathbf e_4$ of the super-loop acts on
$SL(2,\mathbb C)$ by the inversion and on $\mathbb R^{1,3}_{\rm LL}$ by the parity, Eq.~(\ref{eq:inversion}),
and together with $SL(2,\mathbb C)$ it generates a Moufang loop, Eq.~(\ref{eq:chein}). No group contains
such an element (check C-3). Hypothesis that fails: H-grp.

\emph{Counterexample 3: a parallelizable manifold of Lorentz and spinor directions without a group
structure.} The super-loop manifold carries frame fields along the six Lorentz and the eight spinor
directions whose brackets, Eq.~(\ref{eq:frames}), mix the two kinds and depend on the point through the
associator. Its real forms, $S^7_{\rm LL}$ and the pseudo-sphere $S^{3,4}_{\rm LL}$ of the $4+4$ manifold,
have the same property. Hypothesis that fails: H-grp (and H-Lie for the frames).

\emph{Counterexample 4, weaker: spinor charges of the Poincar\'e algebra with a twisted Lorentz action.}
With translations present, no Malcev algebra of our families lets spinor charges close into Lorentz
generators (Theorem M(c)); the only non-Lie extension is the abelian module on which the Lorentz
generators act with the twist $-2$: spinor charges that commute with the translations and with each
other but transform under a Lorentz action that is not a representation (check M-4). Its untwisted
version is a Lie algebra, which Coleman and Mandula exclude as a symmetry of a scattering matrix through
H-op and H-dyn; the twisted version escapes H-Lie as well, but, having no operator realization, it is
not a symmetry of a scattering matrix of an ordinary quantum theory either.

What none of these is: a symmetry of the scattering matrix of a relativistic quantum field theory on a
Hilbert space. Theorem R forbids it. The counterexamples show which conclusions of the no-go theorems
rest on associativity alone: the separation of Lorentz and internal generators (Theorem L versus
Counterexample 1) and the realization of every element by an automorphism (Counterexample 2).
